%&pdftex
% JUAN LUIS VARONA

% CONGRESO UNIMAC'96

% FECHA COMIENZO DOCUMENTO: 15 DE FEBRERO DE 1996
% FECHA \'ULTIMA MODIFICACI\'ON: 6 DE AGOSTO DE 1996
% A\~NADIDA REFERENCIA A LAS ACTAS: 26 DE AGOSTO DE 2004

% ----------------------------------------------------------

\magnification=\magstep1
\pretolerance=2000 % Puesto un n\'umero alto para evitar cortes de palabras
\tolerance=3000
%\overfullrule=0pt
\frenchspacing % Para no dejar doble espacio despu\'es de los puntos
\parskip=4pt plus 3pt minus 2 pt % Un poco m\'as de espacio entre p\'arrafos
% (Jugar con esto para ajustar el n\'umero de p\'aginas)

% \nopagenumbers

\font\titulobf=cmbx12 at 18truept

\font\eightrm=cmr8
\font\eighttt=cmtt9

\font\sc = cmcsc10

% THE ORIGINAL \LaTeX LOGO IS DEFINED HERE:
\def\LaTeX{{\rm L\kern-.36em\raise.3ex\hbox{\sc a}\kern-.15em
    T\kern-.1667em\lower.7ex\hbox{E}\kern-.125emX}}

%%%%%%%%%%%%%%%%%%%%%%%%%%%%%%%%%%%%%%%%%%%%%%%%%%%%%%%%%%%%%%%%%%%%%%%%%%%
%%%%%%%%%%%%%%%%%%%%%%%%%%%%%%%%%% TEXTO %%%%%%%%%%%%%%%%%%%%%%%%%%%%%%%%%%


\centerline{\titulobf HERRAMIENTAS INFORM\'ATICAS PARA}
\medskip
\centerline{\titulobf ESCRIBIR \TeX\ EN CASTELLANO}
\footnote{}{{\sevenrm El trabajo investigador del autor est\'a
subvencionado con la ayuda ULR-96PYB56JVM de la Universidad de La Rioja.}}
\footnote{}{{\sevenrm Esta ponencia fue presentada en el congreso
UNIMAC 96, celebrado del 19 al 21 de septiembre de 1996 en el
Centro Polit\'ecnico Superior de la Universiad de Zaragoza.
Acompa\~nado de las ``herramientas inform\'aticas'' que se citan,
fue publicado en las Actas del Congreso, en el
CD-ROM UNIMAC 96, editado por la Universidad de Zaragoza.}}

\bigskip

\centerline{\bf Juan Luis Varona Malumbres}
\centerline{Dpto.\ de Matem\'aticas y Computaci\'on. Universidad de La Rioja}
\centerline{Calle Luis de Ulloa s/n, 26004 LOGRO\~NO}

\medskip

\centerline{Correo electr\'onico: {\tt jvarona@siur.unirioja.es}}

\bigskip

% ----------------------------------------------------------

{\narrower \eightrm
RESUMEN: 
Presentamos dos herramientas destinadas a ayudar a los usuarios de \TeX\ en
Macintosh, fundamentalmente cuando escriben en castellano. La primera es un
diccionario en espa\~nol y la segunda es un parche destinado a facilitar la
escritura de de documentos \TeX\ en procesadores de texto de Macintosh. Ambas
est\'an, ya en la actualidad, disponibles a la comunidad de usuarios de \TeX\
a trav\'es de Internet por medio de diversos servidores de {\eighttt ftp}
an\'onimo.
\par}

\bigskip

% ----------------------------------------------------------

\beginsection{INTRODUCCI\'ON}

\TeX\ es una mezcla entre procesador de texto y lenguaje de programaci\'on
dedicado a confeccionar documentos escritos de gran calidad tipogr\'afica.
Estos documentos pueden usarse sin cambio alguno en todo tipo de ordenadores
que cuenten con una implementaci\'on de \TeX, y el documento impreso que se
obtiene es independiente del ordenador o impresora usado (l\'ogicamente, las
copias impresas tendr\'an mayor o menor calidad dependiendo de la
resoluci\'on de la impresora). \TeX\ es especialmente \'util cuando se desea
escribir documentos que contienen expresiones matem\'aticas, por complejas que
\'estas sean.

El n\'ucleo de programa \TeX\ fue realizado por D. E. Knuth, y es de libre
disposici\'on para todo tipo de usuarios ---tanto acad\'emicos como los que
desean hacer de \'el un uso comercial---, incluyendo el c\'odigo fuente. En
esto, junto con su gran calidad tipogr\'afica, potencia, y estabilidad frente
a errores, radica su \'exito actual entre las personas que escriben documentos
de tipo cient\'{\i}fico, principalmente matem\'aticos, f\'{\i}sicos e
inform\'aticos. Existen implementaciones de \TeX\ para gran cantidad de
ordenadores, unas gratis y otras de tipo comercial (aunque siempre baratas, ya
que D. E. Knuth no cobra cantidad alguna por usar el n\'ucleo del programa).
Por ejemplo, por citar los m\'as usados, existe para compatibles PC con MS-DOS
o Windows, para equipos UNIX y para Macintosh.

Un fichero fuente de \TeX\ es un archivo en c\'odigo ASCII puro que puede ser
le\'{\i}do por cualquier ordenador, y no se pierde ninguna informaci\'on al
interpretarlo en diversas m\'aquinas o sistemas operativos. Adem\'as, se
transmite por correo electr\'onico sin ning\'un problema, y es la forma
m\'as c\'omoda, segura y f\'acil de compartir documentos complejos a
trav\'es de Internet.

La mayor parte de los temas relacionados con \TeX, as\'{\i} como c\'odigo
fuente e implementaciones para muchos ordenadores, se divulgan y se hacen
disponibles al p\'ublico en general a trav\'es de CTAN (Comprehensive
\TeX\ Archive Network) por medio de ftp an\'onimo en Internet. Actualmente,
la distribuci\'on se hace, oficialmente, a trav\'es de dos servidores, que
esencialmente contienen los mismos archivos. Estos son
$$
\vbox{
\halign{#\hfil&\quad#\hfil\cr
{\it Direcci\'on Internet del ordenador}&
{\it Situaci\'on}\cr
{\tt ftp.uni-stuttgart.edu}&
Universidad de Stuttgart (Alemania)\cr
{\tt ftp.tex.ac.uk}&
Cambridge University (Reino Unido)\cr
}}
$$

En la bibliograf\'{\i}a pueden verse varios libros que constituyen la
referencia obligada de cualquiera que desee profundizar en \TeX\ [3] o en
\LaTeX\ [2, 4]. As\'{\i} mismo, incluimos citas a varios art\'{\i}culos que
informan de la direcci\'on en que se orienta el desarrollo actual basado en
\TeX\ [5] y los intentos para lograr independizarlo lo m\'as posible del
idioma ingl\'es [1].

Referente al uso y desarrollo de \TeX\ para usuarios hispanoablantes, existe un
grupo de usarios (en formaci\'on) denominado Cervan\TeX, que est\'a
coordinado por Jos\'e Ram\'on Portillo, de la Universidad de Sevilla. Puede
consultarse la p\'agina Web del grupo en
{\tt http://gordo.us.es/Actividades/GUTH}. As\'{\i} mismo, existe una lista de
respuestas a preguntas frecuentes sobre asuntos relacionados con el uso de
\TeX\ en castellano; dicha lista est\'a mantenida por Santigo Vila, de la
Universidad de Extremadura, y puede accederse a ella en {\tt
ftp://ftp.unex.es/pub/tex/faq/}. Consultando cualquiera de las dos
direcciones anteriores puede obtenerse todo tipo informaci\'on sobre \TeX\ en
castellano, as\'{\i} como los modos de entrar en contacto con los diferentes
foros de debate que existen.

En lo que a nosotros respecta, aqu\'{\i} presentamos dos herramientas
destinadas a ayudar a los usuarios de \TeX\ en Macintosh, fundamentalmente
cuando escriben en castellano. La primera es un diccionario en espa\~nol y la
segunda es un parche destinado a facilitar la escritura de documentos
\TeX\ en procesadores de texto de Macintosh. Ambas est\'an, ya en la
actualidad, disponibles a la comunidad de usuarios de \TeX\ a trav\'es de
Internet por medio de diversos servidores de {\tt ftp} an\'onimo. Aparte de
esto, el que escribe se ofrece, desde estas l\'{\i}neas, a facilit\'arselas a
cualquiera que se las solicite.

% ----------------------------------------------------------

\beginsection{DICCIONARIO PARA \TeX\ EN CASTELLANO}

Para realizar el diccionario en formato inform\'atico adecuado para \TeX\ no
basta con tener en el ordenador un diccionario com\'un de los que venden en
las librer\'{\i}as. En realidad, cada palabra debe aparecer con todas sus
variaciones permitidas al a\~nadirle distintos prefijos y sufijos. Por
ejemplo, los sustantivos y adjetivos deben aparecer en masculino, femenino,
singular y plural; y los verbos conjugados. Y hay que transformar los
diversos caracteres no ingleses (\~n, acentos) a formato entendible por \TeX.

S\'{\i} que existen programas (verificadores ortogr\'aficos) capaces de
chequear documentos \TeX\ a partir de un diccionario. Por ejemplo,
``Excalibur'' para Macintosh y ``Amspell'' para PC (y no podemos dejar de
citar ``ispell'' que, aunque con un enfoque distinto, est\'a disponible para
equipos UNIX). Pero, l\'ogicamente, no tienen ninguna utilidad para las
personas que escriben en espa\~nol sin un diccionario en ese idioma. Hasta
ahora, hab\'{\i}a disponibles en multitud de idiomas, pero no hab\'{\i}a
ninguno en espa\~nol de dominio p\'ublico y suficientemente completo.

Por otra parte, tambi\'en existen diccionarios comerciales para multitud de
procesadores de textos, por ejemplo, para Microsoft Word o para Word Perfect,
por citar dos de los m\'as usuales. Pero, en primer lugar, hay que comprar una
licencia para poderlos usar, es decir, no son de dominio p\'ublico. Y, en
segundo lugar, un diccionario comercial de alg\'un procesador de textos no
sirve para nada a la hora de chequear documentos \TeX, pues encuentra como
errores todos los comandos de manejo de \TeX. Adem\'as, todos ellos son de
car\'acter general y no est\'an especializados. Las a\~nadidos de palabras que
pueden hacerse funcionan de manera muy poco efectiva. El usuario no tiene
acceso al diccionario fuente y no puede modificarlo a su antojo.

El que escribe ha elaborado un diccionario en espa\~nol que permite verificar
la ortograf\'{\i}a de documentos \TeX, y que contiene casi 250.000 palabras.
El diccionario se maneja a trav\'es de ``Excalibur'', de Robert Gottshall y
Rick Zaccone. Es de destacar que es el \'unico en castellano que, adem\'as de
estar disponible gratis, cuenta con un n\'umero suficiente de palabras
para ser comparable, e incluso mejor, que el de los programas comerciales.
Adem\'as, este diccionario incluye multitud de palabras relacionadas con la
matem\'atica, f\'{\i}sica e inform\'atica, que son las disciplinas donde
m\'as se utiliza \TeX.

El diccionario puede conseguirse, accediendo mediante {\tt ftp anonymous}, en
el directorio {\tt systems/mac/support/excalibur}) de CTAN, o en el directorio
{\tt pub/mac/Excalibur-dictionaries} de {\tt ftp.eg.bucknell.edu}.
Lamentablemente, como la Universidad de La Rioja no tiene ning\'un servidor de
{\tt ftp}, no es posible hacerlo disponible a la comunidad cient\'{\i}fica (y
al p\'ublico en general) en nuestra propia Universidad.

% ----------------------------------------------------------

\beginsection{PARCHE PARA EDITORES DE TEXTO}

Que los archivos fuente de \TeX\ sean c\'odigo ASCII puro requiere varias
complicaciones (al menos a primera vista):

\item{(i)}
En un documento \TeX, por ejemplo, para que, una vez ``compilado'' el
documento, aparezca la letra griega $\alpha$ hay que teclear, dentro de
c\'odigo fuente, la sucesi\'on de caracteres {\tt $\backslash$alpha}.

\item{(ii)}
Las letras no inglesas no est\'an, en principio, definidas, y se obtienen, en
el documento fuente, a trav\'es de distintos artificios. As\'{\i}, una
interrogaci\'on abierta (?`) se escribe como {\tt ?{}`}, una e\~ne como {\tt
$\backslash$\~{}n}, y una a acentuada como {\tt $\backslash$'{}a}.

Estas caracter\'{\i}sticas permiten escribir con \TeX\ documentos de gran
complejidad tipogr\'afica pero, en alguna medida, se puede considerar que
representan una desventaja con respecto a los procesadores de texto
denominados WYSIWYG.

Para intentar solventarla, se han elaborado una serie de parches que se pueden
a\~nadir a un procesador de texto y que facilitar la labor de la persona que
escribe documentos \TeX. Para usarlos no se requiere ning\'un INIT, sino que
los parches se incorporan al editor deseado. Se instalan con ``Text Editor
Patches'', de James W. Walker. De esta forma, se logran escribir f\'ormulas
matem\'aticas de manera m\'as sencilla (por ejemplo, para teclear los
caracteres correspondientes a la letra griega $\alpha$ bastar\'{\i}a pulsar
simult\'aneamente las teclas {\it opci\'on\/} y {\it a\/}) y teclear las letras
acentuadas y otros caracteres no ingleses tal como lo har\'{\i}amos en
cualquier otro procesador de texto, transform\'andose autom\'aticamente en la
notaci\'on requerida por \TeX.

Estos parches est\'an desarrollados teniendo en cuenta el tipo de teclado del
que disponga el Macintosh, ya sea QWERTY (como el espa\~nol y el ingl\'es),
AZERTY (franc\'es y flamenco), QWERTZ (alem\'an y suizo) o QZERTY (italiano).

Por \'ultimo, no podemos pasar por alto que, desde la versi\'on 3.0 de \TeX,
es posible usarlo con un juego completo de caracteres de 8 bits, lo cual
incluye acentos y caracteres no ingleses. Aunque esto tiene el inconveniente
de que estos caracteres no tienen los mismos n\'umeros ASCII en todos los
ordenadores, con lo que surgen problemas al intercambiar ficheros entre
distintas plataformas (hay que decir que tales problemas son f\'acilmente
solucionables). El caso es que unas personas prefieren escribir \TeX\ con un
conjunto de caracteres de 7 bits, y otras con uno completo de 8 bits (lo cual
incluye acentos y caracteres no ingleses). As\'{\i}, hemos incluido en los
parches la posibilidad de desconectarse por etapas. Por ejemplo, se pueden
mantener activos los atajos para escribir letras griegas pero no los
correspondientes a los caracteres no ingleses.

Todos estos parches para \TeX, junto con la aplicaci\'on ``Text Editor
Patches'' que los instala, as\'{\i} como documentaci\'on completa explicando
su funcionamiento, pueden conseguirse mediante {\tt ftp anonymous} en el
directorio {\tt text} de {\tt info-mac} (es decir, en {\tt
sumex-aim.stanford.edu}), el servidor de programas para Macintosh m\'as
conocido y que cuenta con multitud de {\it mirrors\/} de f\'acil acceso.

% ----------------------------------------------------------

\beginsection{Referencias}

\item{[1]} {\sc J. Braams}, {\it Babel, a multilingual style-option system for
use with \LaTeX's standard document styles}, TUGboat {\bf 12} (1991), 291--301.

\item{[2]} {\sc M. Goossens, F. Mittelbach y A. Samarin}, {\it ``The \LaTeX\
Companion''}, Addison Wesley, Reading, 1994.

\item{[3]} {\sc D. E. Knuth}, {\it ``The \TeX book''}, Addison-Wesley, Reading,
1994.

\item{[4]}{\sc  L. Lamport}, {\it ``\LaTeX\ ---A Document Preparation System---
User's Guide and Reference Manual''}, 2\raise1ex\hbox{\underbar{\sevenrm a}}
ed., Addison-Wesley, Reading, 1994.

\item{[5]} {\sc F. Mittebatch y R. Sch\"opf}, {\it Towers \LaTeX3.0}, TUGboat
{\bf 12} (1991), 74--79.

\bye
